\documentclass[12pt,letterpaper]{article}

% Plantilla académica SENA.
% Compilar con LaTeX Workshop o: latexmk -pdf plantilla_trabajo.tex
% Requiere biber y el paquete biblatex-apa.

\usepackage[spanish,es-nodecimaldot,shorthands=off]{babel}
\usepackage[T1]{fontenc}
\usepackage[utf8]{inputenc}
\usepackage{lmodern}
\usepackage{geometry}
\usepackage{setspace}
\usepackage{csquotes}
\usepackage{hyperref}
\usepackage{fancyhdr}
\usepackage{titlesec}
\usepackage[style=apa,backend=biber]{biblatex}

% ==================== DATOS CONFIGURABLES ====================
\newcommand{\institucion}{Servicio Nacional de Aprendizaje -- SENA}
\newcommand{\programa}{[Nombre del programa de formación]}
\newcommand{\tituloTrabajo}{[Título de la actividad]}
\newcommand{\codigoEvidencia}{[Código de la evidencia]}
\newcommand{\aprendiz}{Kevin Andres Granados Daza}
\newcommand{\instructor}{[Nombre del instructor(a)]}
\newcommand{\ficha}{[Número de ficha]}
\newcommand{\centro}{[Nombre del centro de formación]}
\newcommand{\ciudad}{Bogotá D. C.}
\newcommand{\fechaEntrega}{[Fecha de entrega]}
% =============================================================

\geometry{top=2.54cm,bottom=2.54cm,left=2.54cm,right=2.54cm}
\onehalfspacing
\setlength{\parindent}{1.27cm}
\setlength{\parskip}{0pt}
\addbibresource{referencias.bib}
\setlength{\headheight}{15pt}
\pagestyle{fancy}
\fancyhf{}
\fancyhead[R]{\thepage}
\renewcommand{\headrulewidth}{0pt}
\titleformat{\section}{\normalfont\bfseries\centering}{\thesection.}{0.5em}{}

\begin{document}

\begin{titlepage}
  \centering
  \vspace*{2.5cm}
  {\Large \tituloTrabajo\par}
  \vfill
  {\large \aprendiz\par}
  \vspace{1.5cm}
  {\large \institucion\par}
  {\large \programa\par}
  \vspace{1cm}
  {\large Evidencia: \codigoEvidencia\par}
  \vspace{1cm}
  {\large Instructor(a): \instructor\par}
  {\large Ficha: \ficha\par}
  {\large Centro de formación: \centro\par}
  \vfill
  {\large \ciudad\par}
  {\large \fechaEntrega\par}
\end{titlepage}

\section{Introducción}
Escriba el propósito del documento, el contexto de la actividad y los temas que se desarrollarán.

\section{Desarrollo}
Desarrolle la actividad según la guía. Para una cita narrativa use \textcite{ejemploLibro}; para una cita parentética use \parencite{ejemploLibro}.

\section{Conclusiones}
Redacte conclusiones propias que sinteticen los aprendizajes, resultados y aportes del trabajo.

\printbibliography[title={Referencias}]

\end{document}
